% =====================================================================
% RESULTS -- structured by the claim set (PROJECT-RESEARCH-NOTES §33).
% Budget: ~2.5 pages of the 6.  Slots: F0, F1, F2 + tables T3 (+defs).
% Policy rule: DP appears only in the Claim-4 material (R3/R4);
% R1/R2 are argued BL-only.
% =====================================================================
\section{Results}
\label{sec:results}
\subsection{Experimental Setup and Tools}
\label{sec:setup}
All results are obtained with Noxim3D, a cycle-accurate SystemC network-on-chip simulator extended to three dimensions. The substrate is a 6×6×3 mesh of 108 IMC tiles, each holding 128×128 crossbars; links and injection/ejection ports carry one 4-byte flit per 1 ns cycle, giving 4 GB/s per channel, with injection and ejection modelled as independent channels. 
Routing is odd-even-balanced adaptive routing [dahir2013oeb], on top of which we compare two selection policies: a local buffer-occupancy heuristic (BL) and a distributed multi-hop congestion cost-to-go field (DP) \cite{makdpn2009}, whose cost field reconverges every 648 cycles and is clocked at 4× the network clock.

Workloads are ResNet-50, VGG-16 and DeiT-S, converted offline into statistical traffic tables whose rows carry per-pair injection rates and phase windows reproducing the layer-sequential structure of inference; one inference period is 38 536, 115 606 and 58 097 cycles respectively. Each measurement point is warmed up for 1 944 cycles and averaged over three seeds, with offered load swept by a scalar multiplier on all injection rates; unless stated otherwise, comparisons are made inside the knee window, defined as the load range in which mean BL latency stays below 30× its free-flow value. Placement search and all offline metrics (PF, PL, CC, ES) are computed in Python from the traffic tables and the permutation alone. The results below aggregate approximately 36 000 simulation runs.

Traffic is generated from statistical tables rather than dependency-driven traces. This bounds absolute latency accuracy, but not the comparisons the paper makes: every result is a ratio between placements measured under the same model, so a systematic bias cancels unless it reorders them. The model follows established practice for DNN NoC studies \cite{krishnan2022impact}, and the phase windows preserve the layer-sequential structure that the temporal results depend on.



\subsection{Metrics and operating point}
\label{sec:defs}
For a packing $(c,r,s)$ and a placement of its tiles onto the $6{\times}6{\times}3$
mesh, all predictors below are computed offline from the traffic table alone.

\begin{itemize}
  \item \textbf{Port floor} $\PF = \max(\mathrm{PIL},\mathrm{PEL})$: the peak
  per-node injection (PIL) or ejection (PEL) rate, in flits/cycle. Both terms are
  fixed by the flow graph, so $\PF$ is \emph{placement-invariant}: remapping
  relabels which node is hot, never how hot it is (verified identical to six
  decimal places across random permutations).
  \item \textbf{Peak link load} $\PL$: the maximum offered rate on any
  router-to-router channel when each flow spreads evenly over its admissible
  minimal paths under the routing function. $\PL$ is the only term a placement
  moves.
  \item \textbf{Binding load} $\BIND = \max(\PF, \PL)$.
  \item \textbf{Escapable fraction} $\Esc = 1 - \PL_f/\PL$ at the peak link,
  where $\PL_f$ counts only flows for which \emph{every} admissible path crosses
  that link. $\Esc$ is the share of the hot link's load that a selection policy
  could, in principle, divert.
  \item \textbf{Per-mapping capacity} $\kstar$: the injection scale at which a
  mapping reaches twice its own free-flow delay, interpolated from a load
  ladder. Mappings are compared at their own $\kstar$, never at one global
  load~\todo{cite the k* method note}.
\end{itemize}

Load is reported as the fraction $\PF{\cdot}k$ of the binding port's capacity;
the three workloads all knee at 44--48\% of the dominant reduction sink's
ejection bandwidth, so this normalization is workload-independent.

\subsection{Packing sets the floor, and minimum traffic does not minimize it}
\label{sec:r1}

% ---- Claim 0 + Claim 2 (the PO half). Figure F0. BL-only section. ----

At fixed density $c$, the orientation $(r,s)$ decides how much of the partial-sum
reduction is internalised on-tile and how much reaches the NoC, and thereby sets
$\PF$: across ResNet-50's eleven feasible $(c,r,s)$ points $\PF$ spans
$0.258$--$1.494$ flits/cycle ($5.8\times$) before a single placement decision is
made. Placement, by contrast, moves only $\PL$ (span $7.5\times$ within one
packing). The two terms of $\BIND$ therefore belong to different design layers
--- $\PF$ to packing, $\PL$ to placement --- and the rest of this section
establishes which layer holds control in which regime.

Fig.~\ref{fig:designspace} plots every feasible packing of the three workloads
in the (traffic volume, $\PF$) plane. The classical pre-NoC objective ---
minimize communication volume --- and the architecture-aware objective ---
minimize $\PF$ --- select \emph{different} orientations in five of the eight
$(workload, c)$ cells where more than one orientation exists. The disagreement
is not marginal: at $c{=}16$, ResNet-50's minimum-volume orientation $(4,4)$
carries a $2.78\times$ higher floor than $(2,8)$, which costs only $8\%$ more
traffic; VGG-16 at $c{=}32$ shows the same pattern ($2.57\times$ the floor,
$7\%$ traffic). Where the trade is expensive the figure says so: DeiT-S buys
its $1.8$--$3.2\times$ floor reductions with $1.6\times$ the traffic. Minimum
bytes selects the burstiest packings; it is the wrong objective whenever the
floor, not the volume, binds --- which Sec.~\ref{sec:r2} shows is exactly the
regime that matters.

\begin{figure}[t]
  \centering
  \includegraphics[width=\linewidth]{fig0_packing_designspace.pdf}
  \caption{The packing design space: traffic volume against port floor $\PF$,
  all feasible $(c,r,s)$ points, argmins taken \emph{within} each density $c$.
  Arrows run from the minimum-volume orientation to the minimum-$\PF$
  orientation at the same $c$; rings mark cells where the two objectives agree.
  This figure argues about the floor; that the floor determines delay is
  established in Fig.~\ref{fig:hinge}.}
  \label{fig:designspace}
\end{figure}

\subsection{The regime boundary: delay is flat below the floor and steep above}
\label{sec:r2}

% ---- Claims 1-3. Figure F1 (BL only). Hinge + licence + margin hedge. ----

Fig.~\ref{fig:hinge} plots mean delay against $\PL$ for 24 placements of one
packing spanning $\PL = 0.127$--$0.490$ ($3.9\times$) at two loads, with path
variety pinned so $\PL$ is the only moving predictor. Delay does not respond to
$\PL$ until $\PL$ crosses $\PF$, and responds sharply after: a hinge regression
at the floor yields $R^2 = 0.735$ (vs.\ $0.404$ linear), with a below-floor
slope indistinguishable from zero ($t=-0.16$) and an above-floor slope of
$+655$ cycles per unit load ($t=+5.13$, $p<10^{-4}$). The hinge replicates at a
second load level ($R^2 = 0.705$; slope ratio $64$--$72\times$ across
$k$)~\todo{regenerate fig with the E1 grid; cite exact fig file}.

\begin{figure}[t]
  \centering
  \includegraphics[width=\linewidth]{hill_pl_hinge.png}
  \caption{Mean delay against peak link load $\PL$ under local (buffer-level)
  selection, $\PF$ marked. Below the floor a $3.07\times$ spread in $\PL$ moves
  delay $1.15$--$1.18\times$ (n.s.); above it a $1.20\times$ spread moves delay
  $1.71$--$1.85\times$. \todo{replace with E1-grid version, BL only}}
  \label{fig:hinge}
\end{figure}

\textbf{Below the floor, the packing orientation is the determining layer.}
With $\BIND = \PF$, delay tracks the floor itself
%(elasticity ${\approx}4.4$ against $\PF{\cdot}k$)
\todo{E5: PF varied through $(r,s)$ at matched $k$ -- currently varied through
$k$ only; run E5 before submission}, and placement is inert: seventeen
placements spanning $3.07\times$ in $\PL$ move mean delay $1.15\times$
(n.s.). We stress the constructive reading: \emph{verifying $\PL < \PF$ frees
the entire mapping budget} for objectives that are not latency --- thermal
spreading, energy, load balance --- at no delay or p99 penalty. On DeiT-S this
licence replicates end-to-end: at two packings held below their floors, a
placement searched for escapable fraction plus a free choice of selection
policy buys $0.99$--$1.11\times$ over the classical minimum-hop mapping with
local selection (1{,}344 runs; indistinguishable from parity at most loads).

The licence carries one condition: \emph{margin}. Among eight placements all
below the floor, capacity $\kstar$ still varies $1.28\times$ and correlates
with proximity to the boundary ($r(\PL/\PF, \kstar) = -0.70$); at mid-load the
two placements nearest the floor degrade $13.6\times$ earlier than the rest.
The budget is free at $\PL \lesssim 0.6\,\PF$, not asymptotically up to the
boundary.

\textbf{Above the floor, control passes to placement.} Across the 24-placement
grid of Sec.~\ref{sec:r4}, $\PL$ predicts per-mapping capacity at
$r(\PL,\kstar) = -0.60$ ($p=0.002$), and capacity spans $1.50\times$ across
placements of the \emph{same} packing. This is also the first simulator
validation of the offline $\PL$ model ($R^2 = 0.37$): a moderate capacity
predictor, no more~\todo{VGG replication: $r(\PL,\kstar)=-0.72$, BL arm; DP arm
and DeiT-S pending (E7)}.

\subsection{Selection pays only above the floor, and only where load can escape}
\label{sec:r3}

% ---- Claim 4, part 1. Table T3. DP enters the paper HERE. ----

Table~\ref{tab:regime} compares the distributed cost-to-go policy (DP) against
the local buffer-level heuristic (BL) on identical placements, split by regime.
Below the floor the policies are statistically indistinguishable --- no routing
decision reaches a flit that has not entered the network --- while above it DP
removes $7$--$10\%$ of mean delay. The gain is not universal: it requires
escapable load. Pooled across five workload/packing arms ($n{=}71$
placements), DP's capacity gain tracks the escapable fraction $\Esc$ on
ejection-bound packings ($r = +0.59$, $p=0.0002$) and vanishes on
injection-bound ones ($r = +0.07$, n.s.), where the binding port is upstream of
every routing decision.

\begin{table}[t]
  \caption{DP vs.\ BL mean delay on identical placements, by regime
  (24 placements, $n{=}30$ seeds).}
  \label{tab:regime}
  \centering
  \begin{tabular}{lcccc}
    \toprule
    & \multicolumn{2}{c}{$k=1.60$} & \multicolumn{2}{c}{$k=1.80$} \\
    regime & $\Delta$ & $p$ & $\Delta$ & $p$ \\
    \midrule
    below floor ($n{=}17$) & $-0.0\%$ & $0.98$ & $-1.1\%$ & $0.37$ \\
    above floor ($n{=}7$)  & $\mathbf{-10.0\%}$ & $0.0005$ & $\mathbf{-6.7\%}$ & $0.008$ \\
    \bottomrule
  \end{tabular}
\end{table}

\subsection{Adaptive selection returns the mapping freedom that congestion took}
\label{sec:r4}

% ---- Claim 4, part 2. Figure F2 (c4_spread). The variance story. ----

The headline effect of DP is not its mean --- it is its \emph{variance across
placements}. Fig.~\ref{fig:spread} compares the two policies over placements of
equal $\PL$ above the floor. DP does not make the best placements faster (BL
matches or beats it there); it removes the tail: across 21 equal-$\PL$
placements the delay spread contracts from $1.59\times$ to $1.45\times$ (CV
$0.145 \to 0.101$), and in a population containing BL-collapse cases from
$64\times$ to $8.6\times$. On the capacity axis the same grid gives DP a mean
$+10.3\%$ $\kstar$ gain ($p=2.3\times10^{-4}$, positive on 18/24), sized by
$\Esc$: a two-predictor model $\mathrm{gain} \sim \Esc + \PL$ reaches adjusted
$R^2 = 0.55$, and every placement where DP loses is low-$\Esc$.

The design consequence mirrors Sec.~\ref{sec:r2}: once a $\PL$ level is fixed,
the choice \emph{among} equal-$\PL$ placements matters far less under DP than
under BL --- adaptive selection extends the below-floor mapping licence into
the congested regime. What it does not do is substitute for the placement
decision itself: the $\kstar$ spread across $\PL$ levels survives under both
policies.

\begin{figure}[t]
  \centering
  \includegraphics[width=\linewidth]{c4_spread.png}
  \caption{Delay across placements of equal peak link load, BL vs.\ DP, above
  the floor ($k=0.65$). DP compresses the spread across placements rather than
  shifting its floor. \todo{final styling pass; confirm k in caption}}
  \label{fig:spread}
\end{figure}

\subsection{Scope and threats to validity}
\label{sec:r5}

% ---- the caveats that must travel with any citation ----

$\PL$, $\PL_f$ and $\Esc$ are offline path-enumeration models; only $\PL$ has
been validated against the simulator, and only as a moderate capacity predictor
($R^2=0.37$). The even-spreading assumption behind $\PL$ models an ideal
balancer, which BL is not --- BL underperforms its own $\PL$ estimate in
proportion to $\Esc$. All results use one $6{\times}6{\times}3$ mesh; mesh-size
generality is out of scope. Partial sums are carried at INT32, a convention
that \emph{overstates} the reduction traffic the contribution exploits; the
placement rankings survive at 16-bit for ResNet and VGG. Cross-workload
replication is complete for the hinge and the licence, partial for the
selection results~\todo{close with E7 + VGG DP arm, then finalize}.
